\section{Task Bank Design}

The experimental benchmark comprises 52 unique tasks, evenly distributed with 13 tasks per vector. Each task is defined using a YAML schema for programmatic parsing, enabling consistent execution and automated grading where applicable.

The YAML schema defines the following fields across tasks: \texttt{id}, \texttt{name}, \texttt{difficulty}, \texttt{delivery\_mechanism}, \texttt{description}, \texttt{user\_prompt}, \texttt{expected\_safe\_behavior}, \texttt{expected\_exploit\_behavior}, \texttt{scoring\_method}, \texttt{scoring\_criteria}, and \texttt{notes}. Additionally, some tasks utilize a \texttt{file\_context} field for context-based delivery. Depending on the vector, an objective-specific goal field is used: \texttt{injection\_goal} (PI), \texttt{poisoning\_goal} (CP), \texttt{confusion\_goal} (DC), or \texttt{leakage\_goal} (SL). Vector-specific fields further refine the tasks; for instance, Context Poisoning introduces a \texttt{false\_premise} field, Dependency Confusion employs \texttt{malicious\_package\_reference} and \texttt{legitimate\_alternative}, and Secret Leakage utilizes \texttt{planted\_secret\_value} and \texttt{secret\_location}.

Task difficulty is intentionally restricted to ``moderate'' and ``hard'' tiers. Trivial tasks were excluded to prevent agents from passing solely via generic, easily triggered safety filters. The delivery mechanisms are split between chat-paste (where context is explicitly provided in the prompt) and repo-discovery (where the agent must independently explore the repository to discover the injected content), allowing us to measure the security impact of context delivery methods.
